\par\Needspace{6\baselineskip}\medskip
\hypertarget{native-146-article106}{}\pdfbookmark[2]{106.}{article-106}
\begin{center}106.\end{center}
In discussing the most remarkable relations concerning the motion of celestial bodies along conic sections, we must mention the elegant expression for the time in terms of the semimajor axis, the sum $r+r'$, and the chord joining the two positions. For the parabola, this formula appears to have been discovered first \SourceLeaf{147}by the celebrated \textsc{Euler} (Miscell. Berolin. T. VII p. 20), who subsequently neglected it and did not extend it to the ellipse and the hyperbola. Those who attribute the formula to the distinguished \textsc{Lambert} are therefore mistaken, although this geometer deserves credit for independently recovering the forgotten expression and extending it to the other conic sections. Several geometers have already treated the subject, but attentive readers will find the following exposition worthwhile. We begin with elliptic motion.
\par
First of all, note that the angle $2f$ described about the Sun (Art. 88, from which we also take the remaining notation) may be assumed to be less than $360^0$. Clearly, if this angle is increased by 360 degrees, the time increases by the duration of one revolution, or $\dfrac{a^{\frac32}\mathpunct{.}360^0}{k}=a^{\frac32}\times365{,}25$ days. Now if we denote the chord by $\rho$, it is evident that
\[
\rho\rho=(r'\cos v'-r\cos v)^2+(r'\sin v'-r\sin v)^2
\]
\noindent and thus, by equations VIII and IX of Art. 8,
\[
\begin{aligned}
\rho\rho&=aa(\cos E'-\cos E)^2+aa\cos\varphi^2(\sin E'-\sin E)^2\\[.6em]
&=4aa\sin g^2(\sin G^2+\cos\varphi^2\cos G^2)\\[.6em]
&=4aa\sin g^2(1-ee\cos G^2).
\end{aligned}
\]
\noindent Let us introduce an auxiliary angle $h$ such that $\cos h=e\cos G$. To remove all ambiguity, we take $h$ between 0 and $180^0$, so that $\sin h$ is positive. Since $g$ likewise lies between the same limits (if $2g$ reached or exceeded $360^0$, the motion about the Sun would reach or exceed a complete revolution), it follows directly from the preceding equation that
\[
\rho=2a\sin g\sin h,
\]
\noindent provided that the chord is taken to be positive. Since we also have
\[
r+r'=2a(1-e\cos g\cos G)=2a(1-\cos g\cos h),
\]
\noindent it is evident that setting $h-g=\delta$, $h+g=\epsilon$ gives
\[
\begin{aligned}
\mathrm{1)}\qquad&r+r'-\rho=2a(1-\cos\delta)=4a\sin\tfrac12\delta^2\\[.6em]
\mathrm{2)}\qquad&r+r'+\rho=2a(1-\cos\epsilon)=4a\sin\tfrac12\epsilon^2.
\end{aligned}
\]
\par
Finally, we have
\[
kt=a^{\frac32}(2g-2e\sin g\cos G)=a^{\frac32}(2g-2\sin g\cos h),
\]
\noindent\SourceLeaf{148}or
\[
\mathrm{3)}\qquad kt=a^{\frac32}(\epsilon-\sin\epsilon-(\delta-\sin\delta)).
\]
\par
According to equations 1 and 2, the angles $\delta$ and $\epsilon$ can therefore be determined from $r+r'$, $\rho$, and $a$. Using equation 3, the time $t$ can accordingly be determined from the same quantities. If preferred, the formula may be written as follows.
\[
\begin{aligned}
kt=a^{\frac32}\Biggl\{&\operatorname{arc}\cos\frac{2a-(r+r')-\rho}{2a}\\[.6em]
&-\sin\operatorname{arc}\cos\frac{2a-(r+r')-\rho}{2a}\\[.6em]
&-\operatorname{arc}\cos\frac{2a-(r+r')+\rho}{2a}\\[.6em]
&+\sin\operatorname{arc}\cos\frac{2a-(r+r')+\rho}{2a}\Biggr\}\mathpunct{\cdot}
\end{aligned}
\]
\noindent The determination of the angles $\delta$, $\epsilon$ from their cosines leaves an ambiguity that must be examined more closely. Clearly, $\delta$ must lie between $-180^0$ and $+180^0$, and $\epsilon$ between 0 and $360^0$. Even within these limits, each angle seems to admit two determinations, so that the resulting time would admit four. But from equation 5 of Art. 88 we have
\[
\cos f\mathpunct{.}\BareRoot{rr'}=a(\cos g-\cos h)=2a\sin\tfrac12\delta\sin\tfrac12\epsilon:
\]
\noindent Now $\sin\tfrac12\epsilon$ is necessarily positive. We conclude that $\cos f$ and $\sin\tfrac12\delta$ must have the same sign, so that $\delta$ is to be taken between 0 and $180^0$ or between $-180^0$ and 0 according as $\cos f$ is positive or negative, that is, according as the heliocentric motion $2f$ is below or above $180^0$. It is also evident that for $2f=180^0$ we must have $\delta=0$. In this way $\delta$ is fully determined. The determination of the angle $\epsilon$ necessarily remains ambiguous, so that \textit{two} values always result for the time. Deciding which is the true value requires information from elsewhere. The reason is readily seen. Through two given points one can draw \textit{two} different ellipses, both having a focus at the same given point and the same semimajor axis\footnote{Draw a circle of radius $2a-r$ centered at the first position and another of radius $2a-r'$ centered at the second. The other focus of the ellipse evidently lies at an intersection of these circles. Since there are always two intersections in the general case, two different ellipses result.}. Evidently the motion from the first position to the second takes unequal times on these ellipses.
